\documentclass[12pt,a4paper]{article}

% ── Codificación y lenguaje ──────────────────────────────────────────────────
\usepackage[utf8]{inputenc}
\usepackage[T1]{fontenc}
\usepackage[spanish,es-nodecimaldot]{babel}

% ── Tipografía ───────────────────────────────────────────────────────────────
\usepackage{lmodern}
\usepackage{microtype}

% ── Márgenes ─────────────────────────────────────────────────────────────────
\usepackage[top=2.5cm,bottom=2.5cm,left=3cm,right=3cm]{geometry}

% ── Matemáticas ──────────────────────────────────────────────────────────────
\usepackage{amsmath,amssymb}

% ── Figuras y tablas ─────────────────────────────────────────────────────────
\usepackage{graphicx}
\usepackage{booktabs}
\usepackage{array}
\usepackage{float}
\usepackage{caption}
\usepackage{subcaption}
\usepackage{multirow}
\usepackage{diagbox}
\usepackage{colortbl}
\usepackage{rotating}   % para \begin{sidewaysfigure}
\usepackage{pdflscape}  % para landscape solo en esa página

% ── Código fuente ─────────────────────────────────────────────────────────────
\usepackage{listings}
\usepackage[dvipsnames,svgnames,table]{xcolor}
\definecolor{teal}{RGB}{0,128,128}

\definecolor{codeblue}{rgb}{0.13,0.29,0.53}
\definecolor{codegray}{rgb}{0.5,0.5,0.5}
\definecolor{codegreen}{rgb}{0.0,0.5,0.0}
\definecolor{codeback}{rgb}{0.97,0.97,0.97}

\lstset{
  backgroundcolor=\color{codeback},
  basicstyle=\ttfamily\footnotesize,
  keywordstyle=\color{codeblue}\bfseries,
  commentstyle=\color{codegreen},
  stringstyle=\color{red!70!black},
  numberstyle=\tiny\color{codegray},
  numbers=left,
  stepnumber=1,
  frame=single,
  rulecolor=\color{gray!40},
  breaklines=true,
  breakatwhitespace=true,
  tabsize=4,
  language=Python
}

% ── Hyperlinks ────────────────────────────────────────────────────────────────
\usepackage[colorlinks=true,linkcolor=blue!70!black,citecolor=blue!70!black,urlcolor=blue!70!black]{hyperref}

% ── Encabezado/pie de página ──────────────────────────────────────────────────
\usepackage{fancyhdr}
\pagestyle{fancy}
\fancyhf{}
\fancyhead[L]{\small SpectroClass v3.1 — Informe Técnico-Científico}
\fancyhead[R]{\small \thepage}
\fancyfoot[C]{\small UNCuyo · Facultad de Ciencias Exactas y Naturales · 2026}
\renewcommand{\headrulewidth}{0.4pt}
\renewcommand{\footrulewidth}{0.2pt}

% ── Espaciado ─────────────────────────────────────────────────────────────────
\usepackage{setspace}
\onehalfspacing

% ── Bibliografía ──────────────────────────────────────────────────────────────
\usepackage[numbers,sort&compress]{natbib}

% ─────────────────────────────────────────────────────────────────────────────
\begin{document}

% ── PORTADA ──────────────────────────────────────────────────────────────────
\begin{titlepage}
  \centering
  \vspace*{1.5cm}
  {\Large\bfseries Universidad Nacional de Cuyo\\[0.4em]
  Facultad de Ciencias Exactas y Naturales\par}
  \vspace{1.5cm}
  \rule{\linewidth}{1pt}\\[0.8em]
  {\LARGE\bfseries SpectroClass v3.1\par}
  \vspace{0.5em}
  {\large\itshape Sistema Automatizado de Clasificación Espectral Estelar\\
  con Arquitectura Multi-Método\par}
  \rule{\linewidth}{1pt}
  \vspace{1.5cm}

  {\large\bfseries Informe Técnico-Científico\par}
  \vspace{1.2cm}

  \begin{tabular}{rl}
    \textbf{Autor:}  & Roberto Butrón \\[0.3em]
    \textbf{Versión:}  & 3.1 \\[0.3em]
    \textbf{Fecha:}    & Abril 2026 \\[0.3em]
    \textbf{Palabras clave:} & clasificación espectral, sistema MK, \\
                             & anchos equivalentes, redes neuronales, \\
                             & catálogo ELODIE, espectroscopía estelar
  \end{tabular}

  \vfill
  {\small Congreso: Evolución estelar, exoplanetas y dinámica de sistemas estelares\\
  Mendoza, 22--24 de abril de 2026}
\end{titlepage}

% ── RESUMEN ──────────────────────────────────────────────────────────────────
\begin{abstract}
\noindent
Se presenta \textbf{SpectroClass}, un sistema de código abierto para la clasificación
espectral automática de estrellas que integra cuatro métodos independientes mediante
votación ponderada configurable. El sistema procesa espectros en formato TXT y FITS,
normaliza al continuo mediante ajuste local iterativo con rechazo de rayos cósmicos
basado en la desviación absoluta de la mediana (MAD), y mide \textbf{85 líneas
diagnóstico} con ventanas de integración adaptativas. Se calculan \textbf{21 razones
diagnóstico} que cubren la secuencia OBAFGKM completa. La clasificación se realiza
mediante: (1)~un clasificador físico con árbol de decisión de 28 nodos calibrado sobre
estrellas estándar MK, (2)~un árbol de decisión RandomForest entrenado sobre 684
espectros ELODIE, (3)~comparación con plantillas por $\chi^2$ reducido, y
(4)~clasificadores neuronales k-NN y CNN-1D. El sistema de votación ponderada alcanza
una exactitud global del \textbf{86.5\%} sobre el catálogo ELODIE (856 espectros),
superando en $+2.3$ puntos porcentuales al mejor método individual. SpectroClass se
distribuye como aplicación web Flask con interfaz interactiva, exportación CSV/TXT/PDF
y documentación en castellano.
\end{abstract}

\tableofcontents
\newpage

% ════════════════════════════════════════════════════════════════════════════
\section{Introducción}
% ════════════════════════════════════════════════════════════════════════════

\subsection{Contexto y Motivación}

La clasificación espectral estelar es uno de los pilares de la astrofísica moderna.
Desde los trabajos pioneros de Angelo Secchi en el siglo XIX hasta el sistema
Harvard-Draper, la categorización de estrellas por sus características espectrales ha
permitido comprender la física estelar, la evolución química del universo y la
estructura galáctica.

El sistema de clasificación MK (Morgan-Keenan, 1943) ordena las estrellas en siete
tipos espectrales principales, O-B-A-F-G-K-M, por temperatura efectiva decreciente, y
añade clases de luminosidad (I--V) que indican el estado evolutivo.

Con la llegada de \emph{surveys} masivos como SDSS, LAMOST y Gaia, que generan
millones de espectros por ciclo, la automatización de la clasificación se ha vuelto
imprescindible. Los sistemas existentes (MKCLASS, LAMOST pipeline, StarNet)
presentan limitaciones: algunos priorizan la velocidad sobre la trazabilidad física,
otros requieren infraestructura de entrenamiento no disponible localmente, y la mayoría
no incluye documentación en castellano.

\subsection{Problemática}

Los métodos automatizados de clasificación espectral enfrentan cuatro desafíos
principales:
\begin{enumerate}
  \item \textbf{Variabilidad instrumental:} distintas resoluciones y rangos espectrales
        según el espectrógrafo.
  \item \textbf{Contaminación por artefactos:} rayos cósmicos, defectos de CCD y ruido
        fotónico.
  \item \textbf{Efectos físicos:} rotación estelar, velocidad radial y peculiaridades
        químicas modifican los perfiles.
  \item \textbf{Ambigüedad intrínseca:} los límites entre tipos espectrales son
        continuos, no discretos.
\end{enumerate}

\subsection{Objetivos}

\paragraph{Objetivo general.}
Desarrollar un sistema computacional robusto para la clasificación automatizada de
espectros estelares que integre métodos físicos y de aprendizaje automático, con
trazabilidad diagnóstica y distribución de código abierto.

\paragraph{Objetivos específicos.}
\begin{enumerate}
  \item Implementar normalización espectral al continuo con rechazo de rayos cósmicos
        basado en MAD.
  \item Medir anchos equivalentes (EW) para 85 líneas diagnóstico con ventanas
        adaptativas.
  \item Crear un clasificador físico de 28 nodos siguiendo criterios Gray \& Corbally.
  \item Entrenar modelos de ML (RandomForest, k-NN, CNN-1D) sobre el catálogo ELODIE.
  \item Integrar los clasificadores mediante votación ponderada configurable.
  \item Implementar estimación de clase de luminosidad MK (I--V).
  \item Validar el sistema con el catálogo ELODIE (856 espectros).
  \item Distribuir el sistema como aplicación web de código abierto.
\end{enumerate}

% ════════════════════════════════════════════════════════════════════════════
\section{Marco Teórico}
% ════════════════════════════════════════════════════════════════════════════

\subsection{Fundamentos de Espectroscopía Estelar}

\subsubsection{Formación del Espectro}

El espectro estelar resulta de la interacción entre la radiación del continuo
(fotosfera profunda) y la absorción selectiva de capas superiores. En condiciones de
equilibrio termodinámico local (ETL), la función fuente $S_\lambda$ equivale a la
función de Planck $B_\lambda(T)$, y la intensidad emergente se rige por la ecuación de
transferencia radiativa:
\begin{equation}
  \frac{dI_\lambda}{d\tau_\lambda} = I_\lambda - B_\lambda(T)
\end{equation}

Los perfiles de líneas resultan de la convolución de ensanchamientos natural
(Lorentziano), Doppler térmico, de presión (Stark, Van der Waals) y rotacional,
aproximados habitualmente por un perfil de Voigt.

\subsubsection{Ancho Equivalente}

El ancho equivalente (EW) es la medida fundamental para la clasificación cuantitativa:
\begin{equation}
  \mathrm{EW} = \int \left(1 - \frac{F_\lambda}{F_c}\right) d\lambda \quad [\text{\AA}]
\end{equation}
donde $F_\lambda$ es el flujo observado y $F_c$ el flujo del continuo. El EW es
independiente de la resolución instrumental (a diferencia del FWHM) y refleja la
abundancia del elemento y las condiciones físicas de la atmósfera estelar.

\subsection{Sistema de Clasificación Espectral MK}

La Tabla~\ref{tab:mk} resume los siete tipos principales con sus propiedades físicas y
las características espectrales diagnóstico.

\begin{table}[H]
  \centering
  \caption{Secuencia de tipos espectrales MK.}
  \label{tab:mk}
  \renewcommand{\arraystretch}{1.25}
  \begin{tabular}{cllp{6cm}}
    \toprule
    \textbf{Tipo} & \textbf{$T_\mathrm{ef}$ (K)} & \textbf{Color} &
    \textbf{Características principales} \\
    \midrule
    O & $>30\,000$ & Azul & He~\textsc{ii} fuerte, He~\textsc{i} presente, H débil \\
    B & $10\,000$--$30\,000$ & Azul-blanco & He~\textsc{i} máximo en B2, H creciente \\
    A & $7\,500$--$10\,000$ & Blanco & H máximo (Balmer), metales débiles \\
    F & $6\,000$--$7\,500$ & Amarillo-blanco & H decreciente, Ca~\textsc{ii}~K creciente \\
    G & $5\,200$--$6\,000$ & Amarillo & Ca~\textsc{ii} H y K fuertes, banda G (CH) \\
    K & $3\,700$--$5\,200$ & Naranja & Metales dominantes, TiO incipiente \\
    M & $<3\,700$ & Rojo & Bandas moleculares TiO, VO \\
    \bottomrule
  \end{tabular}
\end{table}

Las clases de luminosidad (Tabla~\ref{tab:lum}) codifican el estado evolutivo a través
de indicadores sensibles a la gravedad superficial $\log g$.

\begin{table}[H]
  \centering
  \caption{Clases de luminosidad MK.}
  \label{tab:lum}
  \renewcommand{\arraystretch}{1.2}
  \begin{tabular}{cll}
    \toprule
    \textbf{Clase} & \textbf{Denominación} & \textbf{Indicador diagnóstico principal} \\
    \midrule
    Ia, Ib & Supergigantes & Líneas estrechas, ratios de ionización alterados \\
    II & Gigantes brillantes & Sr~\textsc{ii} / Fe~\textsc{i} elevado \\
    III & Gigantes & Ba~\textsc{ii} moderado \\
    IV & Subgigantes & Ca~\textsc{i} / Fe~\textsc{i} intermedio \\
    V & Enanas (S.P.) & Líneas más anchas (mayor presión) \\
    \bottomrule
  \end{tabular}
\end{table}

% ════════════════════════════════════════════════════════════════════════════
\section{Metodología}
% ════════════════════════════════════════════════════════════════════════════

\subsection{Pre-procesamiento Espectral}

\subsubsection{Rechazo de Rayos Cósmicos — Sigma-Clipping MAD}

El algoritmo de sigma-clipping utiliza la Desviación Absoluta de la Mediana (MAD)
en lugar de la desviación estándar clásica, lo que lo hace resistente a distribuciones
asimétricas propias de los rayos cósmicos en espectros CCD.
La desviación estándar robusta se calcula como $\hat{\sigma} = 1.4826 \cdot \mathrm{MAD}$,
donde el factor $1.4826$ garantiza consistencia con la distribución Gaussiana.
El proceso itera hasta cinco veces o hasta convergencia.

\subsubsection{Normalización al Continuo}

La normalización correcta es crítica para el cálculo de EW. El algoritmo implementa un
enfoque de \emph{dos pasos}:
\begin{enumerate}
  \item \textbf{Primer pase — continuo rápido:} selección del percentil P90 del flujo
        en ventanas de 50 píxeles, seguida de ajuste con spline cúbico.
  \item \textbf{Segundo pase — rechazo relativo:} los rayos cósmicos se rechazan
        comparando el flujo con el continuo local (no el global), lo que evita
        distorsiones en espectros con fuerte gradiente de color.
\end{enumerate}
El resultado es $F_\mathrm{norm}(\lambda) = F_\mathrm{obs}(\lambda) /
F_\mathrm{continuo}(\lambda)$, con el continuo igual a la unidad.
Se destaca que la normalización min-max es \emph{incorrecta} para espectroscopía, ya
que un rayo cósmico en el máximo comprime todo el espectro hacia cero e impide el
cálculo de EW.

\subsection{Medición de Anchos Equivalentes}

La función \texttt{measure\_equivalent\_width()} retorna una tupla de cuatro elementos:
$(\mathrm{EW},\;\mathrm{depth},\;\mathrm{quality\_flag},\;\lambda_\mathrm{med})$.
El \texttt{quality\_flag} codifica tres niveles de calidad (0 = válido, 1 = warning,
2 = descartado) basados en la profundidad mínima de la línea, la posición del mínimo
respecto del centro teórico (tolerancia adaptativa) y la relación S/N local
($\mathrm{SNR} \geq 30$).

Las ventanas de integración se definen por línea en el diccionario \texttt{LINE\_WINDOWS}
(6--30~Å). Para líneas metálicas angostas se utiliza una ventana adaptativa:
$w \approx 4 \times \mathrm{FWHM}$, estimado mediante la función \texttt{estimate\_fwhm()}.
Las 85 líneas diagnóstico implementadas se distribuyen como se muestra en la
Tabla~\ref{tab:lineas}.

\begin{table}[H]
  \centering
  \caption{Distribución de las 85 líneas diagnóstico por especie.}
  \label{tab:lineas}
  \renewcommand{\arraystretch}{1.2}
  \begin{tabular}{lrl}
    \toprule
    \textbf{Especie / Serie} & \textbf{N.º líneas} & \textbf{Uso diagnóstico} \\
    \midrule
    Balmer (H\,\textsc{i}) & 5 & Temperatura A--F, luminosidad \\
    Helio (He\,\textsc{i}) & 6 & Tipos B \\
    Helio (He\,\textsc{ii}) & 4 & Tipos O \\
    N, C, O ionizados & 13 & Subtipos O y B \\
    Silicio (Si) & 8 & Subtipos B \\
    Calcio (Ca) & 6 & F--G--K primario \\
    Magnesio (Mg) & 3 & B tardío -- A \\
    Hierro (Fe) & 13 & G--K--M \\
    Cr, Sr, Ba, Na & 5 & Luminosidad, K--M \\
    Ti, Al & 2 & G--K \\
    CH banda G & 1 & F/G \\
    Bandas TiO & 6 & M \\
    Bandas VO y CaH & 4 & M tardías \\
    Otros metales & 9 & F--K \\
    \midrule
    \textbf{Total} & \textbf{85} & \\
    \bottomrule
  \end{tabular}
\end{table}

\subsection{Razones Diagnóstico}

La función \texttt{compute\_spectral\_ratios()} calcula 21 razones que cubren la
secuencia completa:
\begin{itemize}
  \item \textbf{Tipos O:} He\,\textsc{i}~4471 / He\,\textsc{ii}~4541,
        N\,\textsc{iv}~4058 / N\,\textsc{iii}~4641, Si\,\textsc{iv}~4089 /
        He\,\textsc{i}~4026, N\,\textsc{v}/N\,\textsc{iii},
        O\,\textsc{iii}~5592 / He\,\textsc{i}~5876.
  \item \textbf{Tipos B:} Si\,\textsc{iii}~4553 / Si\,\textsc{ii}~4128,
        Si\,\textsc{iv}~4116 / Si\,\textsc{iii}~4553.
  \item \textbf{Tipos F--G:} Ca\,\textsc{ii}~K / H$\epsilon$,
        Ca\,\textsc{i}~4227 / H$\delta$, Sr\,\textsc{ii}~4077 / Fe\,\textsc{i}.
  \item \textbf{Tipos G--K:} Cr\,\textsc{i}~4254 / Fe\,\textsc{i},
        Na\,\textsc{i}~D / Fe\,\textsc{i}, Y\,\textsc{ii}~4376 / Fe\,\textsc{i}.
  \item \textbf{Tipos K--M:} TiO~4762 / Ca\,\textsc{i},
        TiO / CaH, VO~7434 / CaH.
\end{itemize}

\subsection{Clasificadores}

\subsubsection{Clasificador Físico — Árbol de 28 Nodos}

El clasificador físico implementa los criterios de Gray \& Corbally (2009) en un árbol
de decisión jerárquico ampliado a 28 nodos.
La Figura~\ref{fig:arbol} muestra el diagrama completo del árbol con la codificación
de color por tipo espectral y los 10 nodos nuevos de la versión~3.1 destacados en verde
(\textit{teal}). Los 10 nodos añadidos en v3.1 cubren:
subtipos O3--O4 (N\,\textsc{v}~4604 + O\,\textsc{v}~5114), O tardíos
(N\,\textsc{iii}~4634--4641), B con O\,\textsc{ii}, F5--F9 (Ca\,\textsc{i}~4226/H$\delta$
+ Sr\,\textsc{ii}~4077), F/G con banda CH G, G con Fe\,\textsc{i}~4144/H$\delta$ +
Cr\,\textsc{i}/Fe\,\textsc{i} + Y\,\textsc{ii}~4376, G con tripleta Mg~$b$, K con
triplete Cr\,\textsc{i} (4254/4275/4290~Å) como termómetro primario, K con
Na\,\textsc{i}~D, y M tardías con VO y CaH.

% ── Figura: árbol de decisión ────────────────────────────────────────────────
\begin{landscape}
\begin{figure}[p]
  \centering
  % Para incluir la figura compilada como PDF externo (recomendado):
  %   1. Compilar: pdflatex arbol_decision_28nodos.tex
  %   2. Descomentar la línea \includegraphics y comentar el \fbox
  %
  % \includegraphics[width=0.97\linewidth,keepaspectratio]{arbol_decision_28nodos.pdf}
  %
  % Mientras no se haya compilado el archivo externo, se muestra este recuadro
  % con instrucciones:
  \fbox{\parbox{0.94\linewidth}{\centering\vspace{1cm}
    \textbf{Árbol de Decisión — 28 Nodos (SpectroClass v3.1)}\\[0.5cm]
    \textit{Para mostrar el diagrama TikZ, compilar primero el archivo externo:}\\[0.3cm]
    \texttt{pdflatex arbol\_decision\_28nodos.tex}\\[0.3cm]
    y luego descomentar la línea \texttt{\textbackslash includegraphics} en este archivo.\\[0.3cm]
    \textit{Alternativamente, este archivo puede compilarse de forma independiente\\
    para generar la figura como un PDF standalone.}
    \vspace{1cm}
  }}
  \caption[Árbol de decisión de 28 nodos — SpectroClass v3.1]{%
    Árbol de decisión del clasificador físico de SpectroClass v3.1 (28 nodos,
    ampliado desde 18 en la versión~3.0).
    Los nodos con borde \textcolor{teal}{\textit{teal}} y relleno verde claro son los
    10 nodos nuevos incorporados en v3.1.
    El código de colores refleja el tipo espectral de destino:
    azul~(O), cian~(B), amarillo-naranja~(A), naranja~(F),
    verde~(G), naranja-rojo~(K) y rojo~(M).
    Cada nodo de decisión (\textit{diamante}) cita las líneas espectrales y umbrales
    de EW utilizados; cada nodo terminal (\textit{rectángulo redondeado}) indica el
    tipo MK asignado y el criterio diagnóstico que lo determina.
    Al pie se incluye el módulo de clase de luminosidad (I--V).
  }
  \label{fig:arbol}
\end{figure}
\end{landscape}

\subsubsection{Árbol de Decisión RandomForest}

Se entrena un clasificador RandomForest sobre 684 espectros ELODIE con 24 features de
entrada (17 EW + 6 razones + H\_avg), profundidad máxima 9, balanceo de clases activo
y validación cruzada 5-fold.

\subsubsection{Clasificadores Neuronales}

\paragraph{k-NN.} Opera en el espacio de features escaladas (StandardScaler),
$k=5$ vecinos con ponderación por distancia euclidiana.

\paragraph{CNN-1D.} Red convolucional sobre el espectro remuestreado a grilla uniforme,
con dos bloques Conv1D--MaxPooling--Dropout (32 y 64 filtros), dos capas densas
(128 unidades + Dropout 0.3) y salida Softmax. Optimizador Adam, tasa de aprendizaje
$10^{-3}$.

\subsubsection{Template Matching}

Comparación por $\chi^2$ reducido con plantillas de EW para cada tipo MK:
\begin{equation}
  \chi^2_\nu = \frac{1}{N}\sum_{i=1}^{N}
  \left(\frac{\mathrm{EW}_i^\mathrm{obs} - \mathrm{EW}_i^\mathrm{plantilla}}
  {\sigma_i}\right)^2
\end{equation}
donde $\sigma_i = \max(0.1 \cdot \mathrm{EW}_i^\mathrm{plantilla},\; 0.05\ \text{\AA})$.

\subsection{Votación Ponderada}

La clasificación final combina los cuatro métodos con pesos:
\begin{align}
  w_\mathrm{físico} &= 0.10, &
  w_\mathrm{DT} &= 0.40, &
  w_\mathrm{template} &= 0.10, &
  w_\mathrm{neural} &= 0.40.
\end{align}
El tipo final es $\hat{T} = \arg\max_T \sum_m w_m \cdot p_m(T)$,
donde $p_m(T)$ es la confianza del método $m$ para el tipo $T$.
La confianza global incorpora un factor de consenso que penaliza resultados discordantes.

\subsection{Estimación de Clase de Luminosidad}

El módulo \texttt{luminosity\_classification.py} estima la clase de luminosidad I--V
mediante indicadores sensibles a $\log g$:
Sr\,\textsc{ii}/Fe\,\textsc{i} (gigantes y supergigantes tardías),
Ba\,\textsc{ii}/Fe\,\textsc{i} (gigantes G--K),
Ca\,\textsc{i}/Fe\,\textsc{i} (subgigantes),
TiO/CaH (tipos M) y criterios Stark en H$\beta$ para tipos A--B.
El resultado se combina con el tipo espectral para construir la clasificación MK
completa (p.~ej., G2V, K3\,\textsc{iii}).

% ════════════════════════════════════════════════════════════════════════════
\section{Implementación Computacional}
% ════════════════════════════════════════════════════════════════════════════

\subsection{Arquitectura del Sistema}

SpectroClass sigue una arquitectura en capas (Figura~\ref{fig:arch}):
(1)~carga de datos (TXT / FITS con reconstrucción WCS), (2)~pre-procesamiento
(normalización + EW), (3)~capa de clasificadores, (4)~votación ponderada y
(5)~interfaz web Flask con visualización interactiva.

\begin{figure}[H]
  \centering
  \fbox{\parbox{0.88\linewidth}{\centering\ttfamily\small
    Interfaz Web (Flask)\\
    \rule{0.7\linewidth}{0.4pt}\\
    SpectralValidator — Votación Multi-Método\\
    \rule{0.7\linewidth}{0.4pt}\\
    Físico \quad | \quad RandomForest \quad | \quad Template \quad | \quad Neural\\
    \rule{0.7\linewidth}{0.4pt}\\
    Pre-procesamiento: Normalización + EW + Razones\\
    \rule{0.7\linewidth}{0.4pt}\\
    Carga de Datos (.txt / .fits)
  }}
  \caption{Arquitectura de capas de SpectroClass v3.1.}
  \label{fig:arch}
\end{figure}

\subsection{Módulos Principales}

\begin{description}
  \item[\texttt{spectral\_classification\_corrected.py}]
    Módulo central: \texttt{sigma\_clip()}, \texttt{normalize\_to\_continuum()},
    \texttt{measure\_equivalent\_width()}, \texttt{measure\_diagnostic\_lines()}
    (85 líneas), \texttt{compute\_spectral\_ratios()} (21 razones),
    \texttt{classify\_star\_corrected()} (árbol 28 nodos),
    \texttt{check\_radial\_velocity()}, \texttt{write\_diagnostic\_file()}.
  \item[\texttt{spectral\_validation.py}]
    \texttt{SpectralValidator}: gestiona el ciclo de votación ponderada e integra
    los cuatro clasificadores.
  \item[\texttt{neural\_classifiers.py}]
    \texttt{KNNClassifier}, \texttt{CNNClassifier} y
    \texttt{NeuralClassifierManager} para entrenamiento e inferencia.
  \item[\texttt{luminosity\_classification.py}]
    \texttt{estimate\_luminosity\_class()} y
    \texttt{combine\_spectral\_and\_luminosity()}.
  \item[\texttt{webapp/app.py}]
    Servidor Flask con endpoints: \texttt{POST /upload} (espectro individual),
    \texttt{POST /batch\_process} (lotes), \texttt{POST /train\_neural},
    \texttt{GET /neural\_metrics} y \texttt{GET /health}.
\end{description}

\subsection{Dependencias}

NumPy, SciPy, Astropy (lectura FITS), scikit-learn (RandomForest, k-NN),
TensorFlow/Keras (CNN-1D), Flask (interfaz web), Matplotlib y Pandas.
Instalación mediante \texttt{pip install -r requirements.txt}.

% ════════════════════════════════════════════════════════════════════════════
\section{Validación y Resultados}
% ════════════════════════════════════════════════════════════════════════════

\subsection{Conjunto de Datos — Catálogo ELODIE}

El catálogo ELODIE del Observatorio de Haute-Provence proporciona 856 espectros
de alta calidad (Tabla~\ref{tab:elodie}).

\begin{table}[H]
  \centering
  \caption{Características del catálogo ELODIE utilizado para validación.}
  \label{tab:elodie}
  \renewcommand{\arraystretch}{1.2}
  \begin{tabular}{ll}
    \toprule
    \textbf{Parámetro} & \textbf{Valor} \\
    \midrule
    Total de espectros & 856 \\
    Rango espectral & 3900--6800~\AA \\
    Resolución & $R \sim 42\,000$ \\
    S/N típico & $> 100$ \\
    División entrenamiento / test & 80\,\% / 20\,\% (684 / 171) \\
    Validación cruzada & 5-fold estratificada \\
    \bottomrule
  \end{tabular}
\end{table}

La distribución de la muestra es: F 41.1\,\%, K 30.0\,\%, A 16.2\,\%, B 7.5\,\%,
M 3.2\,\%, O 1.8\,\% y G 0.1\,\%.

\subsection{Métricas de Rendimiento}

La Tabla~\ref{tab:metricas} compara la exactitud de cada método y del sistema combinado.

\begin{table}[H]
  \centering
  \caption{Exactitud global por método (conjunto de test, $N=171$).}
  \label{tab:metricas}
  \renewcommand{\arraystretch}{1.2}
  \begin{tabular}{lccc}
    \toprule
    \textbf{Método} & \textbf{Exactitud} & \textbf{CV 5-fold} \\
    \midrule
    Árbol de decisión (RandomForest) & 84.2\,\% & $82.5 \pm 3.1$\,\% \\
    k-NN ($k=5$) & 81.3\,\% & $79.8 \pm 2.8$\,\% \\
    Clasificador físico & \multicolumn{2}{c}{trazabilidad completa} \\
    Template matching & \multicolumn{2}{c}{referencia diagnóstica} \\
    \midrule
    \textbf{Votación ponderada (todos)} & \textbf{86.5\,\%} & — \\
    \bottomrule
  \end{tabular}
\end{table}

La Tabla~\ref{tab:report} muestra las métricas detalladas del RandomForest por tipo.

\begin{table}[H]
  \centering
  \caption{Reporte de clasificación del árbol RandomForest por tipo espectral.}
  \label{tab:report}
  \renewcommand{\arraystretch}{1.2}
  \begin{tabular}{lcccc}
    \toprule
    \textbf{Tipo} & \textbf{Precisión} & \textbf{Recall} & \textbf{F1} &
    \textbf{Soporte} \\
    \midrule
    A & 0.82 & 0.86 & 0.84 & 28 \\
    B & 0.75 & 0.69 & 0.72 & 13 \\
    F & 0.89 & 0.90 & 0.89 & 70 \\
    K & 0.82 & 0.80 & 0.81 & 51 \\
    M & 0.83 & 0.91 & 0.87 & 6 \\
    O & 1.00 & 0.67 & 0.80 & 3 \\
    \midrule
    \textbf{Promedio ponderado} & 0.84 & 0.84 & 0.84 & 171 \\
    \bottomrule
  \end{tabular}
\end{table}

\subsection{Matriz de Confusión}

La Tabla~\ref{tab:confusion} muestra la matriz de confusión del sistema de votación
ponderada sobre el conjunto de test.

\begin{table}[H]
  \centering
  \caption{Matriz de confusión — sistema de votación ponderada ($N=171$).}
  \label{tab:confusion}
  \renewcommand{\arraystretch}{1.3}
  \begin{tabular}{c|ccccccc}
    \toprule
    \diagbox{Real}{Pred.} & \textbf{O} & \textbf{B} & \textbf{A} &
    \textbf{F} & \textbf{G} & \textbf{K} & \textbf{M} \\
    \midrule
    O & \textbf{2} & 1 & 0 & 0 & 0 & 0 & 0 \\
    B & 0 & \textbf{9} & 3 & 1 & 0 & 0 & 0 \\
    A & 0 & 2 & \textbf{24} & 2 & 0 & 0 & 0 \\
    F & 0 & 0 & 2 & \textbf{63} & 0 & 5 & 0 \\
    K & 0 & 0 & 0 & 5 & 0 & \textbf{46} & 0 \\
    M & 0 & 0 & 0 & 0 & 0 & 1 & \textbf{5} \\
    \bottomrule
  \end{tabular}
\end{table}

Las principales confusiones son B$\leftrightarrow$A (gradiente continuo de temperatura)
y F$\leftrightarrow$K (solapamiento de características metálicas).
El bajo recall de O se debe al tamaño reducido de la muestra (3 espectros de test).

\subsection{Comparación con Sistemas Previos}

\begin{table}[H]
  \centering
  \caption{Comparación con sistemas de clasificación espectral automatizada.}
  \label{tab:comparacion}
  \renewcommand{\arraystretch}{1.2}
  \begin{tabular}{llll}
    \toprule
    \textbf{Sistema} & \textbf{Exactitud} & \textbf{Método} & \textbf{Referencia} \\
    \midrule
    SpectroClass v3.1 & \textbf{86.5\,\%} & Multi-método & Este trabajo \\
    MKCLASS & $\sim$85\,\% & Reglas físicas & Gray \& Corbally (2014) \\
    LAMOST pipeline & 83\,\% & Template matching & Luo et al. (2015) \\
    StarNet & 89\,\% & CNN & Fabbro et al. (2018) \\
    \bottomrule
  \end{tabular}
\end{table}

% ════════════════════════════════════════════════════════════════════════════
\section{Discusión}
% ════════════════════════════════════════════════════════════════════════════

\subsection{Ventajas del Enfoque Multi-Método}

La combinación de clasificadores mejora la exactitud en $+2.3$ puntos sobre el mejor
método individual por tres razones complementarias: (1)~el clasificador físico
proporciona trazabilidad diagnóstica completa (cada decisión cita la línea y el umbral
que la motivó); (2)~el RandomForest captura relaciones no lineales entre EW que los
criterios físicos simplifican; (3)~la CNN-1D opera sobre el espectro completo y
detecta patrones sutiles en el perfil que los EW individuales no capturan.

\subsection{Limitaciones}

\begin{enumerate}
  \item \textbf{Rango espectral:} Si Ca\,\textsc{ii}~K (3933~Å) queda fuera del
        rango del instrumento, el árbol interactivo activa nodos alternativos; sin
        embargo, la confianza en clasificación F--G--K disminuye.
  \item \textbf{Desbalance de clases:} Los tipos O (1.8\,\%) y G (0.1\,\%) del
        catálogo ELODIE son insuficientes para estimar métricas robustas.
  \item \textbf{Velocidad radial:} Sin corrección de VR, las líneas pueden desplazarse
        hasta $\pm 2$~Å, cubierto por la tolerancia adaptativa pero con riesgo de
        confusión en líneas cercanas.
  \item \textbf{Estrellas peculiares:} Ap, Bp, Am y Be requieren criterios adicionales
        no implementados en la versión actual.
\end{enumerate}

\subsection{Mejoras Implementadas en v3.1}

Con respecto a la versión 3.0, la versión 3.1 incorpora: (1)~sigma-clipping MAD,
(2)~normalización de dos pasos relativa al continuo local, (3)~tupla EW de 4 elementos
con \texttt{quality\_flag}, (4)~diccionario \texttt{LINE\_TYPES} por categoría,
(5)~\texttt{estimate\_fwhm()} y ventanas adaptativas, (6)~\texttt{compute\_snr()} con
constantes \texttt{SNR\_MINIMUM = 30} y \texttt{RESOLUTION\_MINIMUM = 2000},
(7)~registro de \texttt{lambda\_measured} y \texttt{window\_width} en diagnósticos,
(8)~6 razones adicionales para tipos O--B, (9)~\texttt{write\_diagnostic\_file()}
genera \texttt{line\_measurements.txt}, y (10)~\texttt{check\_radial\_velocity()}.

% ════════════════════════════════════════════════════════════════════════════
\section{Conclusiones}
% ════════════════════════════════════════════════════════════════════════════

\begin{enumerate}
  \item Se desarrolló SpectroClass v3.1, un sistema de clasificación espectral estelar
        automatizada de código abierto que integra cuatro métodos complementarios con
        trazabilidad diagnóstica completa.
  \item El sistema alcanza una exactitud del \textbf{86.5\,\%} en la secuencia
        OBAFGKM sobre el catálogo ELODIE (856 espectros), superando en $+2.3$~pp al
        mejor método individual.
  \item La normalización MAD de dos pasos y las ventanas adaptativas por FWHM reducen
        los errores sistemáticos en espectros de alta resolución y en tipos O--B
        tempranos.
  \item La arquitectura multi-método combina trazabilidad física (28 nodos Gray \&
        Corbally), aprendizaje automático (RandomForest, k-NN, CNN-1D) y comparación
        con plantillas, con pesos ajustables por el usuario.
  \item La estimación de clase de luminosidad I--V mediante indicadores sensibles a
        $\log g$ permite reportar la clasificación MK completa (p.~ej., K3\,III).
  \item El sistema se distribuye como aplicación web Flask con interfaz interactiva,
        árbol de clasificación visual, exportación CSV/TXT/PDF y documentación completa
        en castellano.
  \item Trabajo futuro: ampliar el catálogo de entrenamiento (SDSS, LAMOST), incorporar
        corrección automática de VR, extender a estrellas peculiares (Ap, Am, Be) y
        añadir incertidumbres bayesianas en la confianza de clasificación.
\end{enumerate}

% ════════════════════════════════════════════════════════════════════════════
\section*{Agradecimientos}
% ════════════════════════════════════════════════════════════════════════════

El autor agradece al Observatorio de Haute-Provence por el acceso público al catálogo
ELODIE, y a la comunidad de Python científico (Astropy, scikit-learn, TensorFlow) cuyas
bibliotecas de código abierto hacen posible este trabajo.

% ════════════════════════════════════════════════════════════════════════════
\bibliographystyle{unsrtnat}
\begin{thebibliography}{9}

\bibitem{gray2009}
  Gray, R.~O. \& Corbally, C.~J. (2009).
  \textit{Stellar Spectral Classification}.
  Princeton University Press.

\bibitem{mkclass2014}
  Gray, R.~O. \& Corbally, C.~J. (2014).
  MKCLASS: A Program for the Automated Classification of Stellar Spectra.
  \textit{AJ}, 147, 80.

\bibitem{luo2015}
  Luo, A.-L. et al. (2015).
  The First Data Release (DR1) of the LAMOST regular survey.
  \textit{RAA}, 15, 1095.

\bibitem{fabbro2018}
  Fabbro, S. et al. (2018).
  An Application of Deep Learning in the Analysis of Stellar Spectra.
  \textit{MNRAS}, 475, 2978.

\bibitem{moultaka2004}
  Moultaka, J. et al. (2004).
  The ELODIE Archive.
  \textit{PASP}, 116, 693.

\bibitem{morgan1943}
  Morgan, W.~W., Keenan, P.~C. \& Kellman, E. (1943).
  \textit{An Atlas of Stellar Spectra}.
  University of Chicago Press.

\end{thebibliography}

\end{document}
